\section{Algorithm H in brief}
\label{sec:background}

This section recalls the parts of Section~7.5.1 that matter for worst
cases. Readers who know pages~1--5 of the pre-fascicle well can skim it,
but should look at two things: the definition of an \emph{input} in
Section~\ref{sec:input}, which fixes exactly what freedom a worst-case
construction has, and the worked round in Section~\ref{sec:worked}.

\subsection{Matchings and augmenting paths}

A bipartite graph has $M$ \emph{girls} and $N$ \emph{boys}, joined by $E$
edges; each edge joins a girl to a boy. A \emph{matching} is a set of
edges with no common vertex. Its edges are called \emph{couples}, and a
vertex that belongs to no couple is \emph{free}. A matching is
\emph{maximal} if no edge can be added to it, and \emph{maximum} if no
matching has more edges. Throughout, $s$ denotes the size of a maximum
matching. We often write $V=M+N$ for the number of vertices.

We follow the pre-fascicle's drawing conventions. When two matchings are
compared, an edge of the reference matching is written $x\sg y$ and an
edge of the current matching $M'$ is written $x\db y$. In particular an
\emph{augmenting path} for $M'$ is a path
\begin{equation}
  g_0\sg b_1\db g_1\sg b_2\db\cdots\db g_{L-1}\sg b_L
  \label{eq:augpath}
\end{equation}
that starts at a free girl, ends at a free boy, and alternates between
edges outside and inside $M'$. It has $2L-1$ edges, and we say it has
\emph{level}~$L$. Exchanging the two kinds of edges along the path
(Fig.~\ref{fig:augpath}) produces a matching with one more couple.
In~\eqref{eq:augpath} no reference matching is involved, and the single
lines mean ``an edge not in~$M'$.'' We describe augmenting paths by
their level; when the number of edges matters we say so. An edge named
without regard to any matching is written by juxtaposition, as in
$g_1b_6$.

\begin{figure}[ht]
\centering
\begin{tikzpicture}[x=1.05cm]
\foreach \row/\y/\name in {0/0/before,1/-1.25/after}{
  \node[girl,free] (g0\row) at (0,\y) {$g_0$};
  \node[boy] (b1\row) at (1,\y) {$b_1$};
  \node[girl] (g1\row) at (2,\y) {$g_1$};
  \node[boy] (b2\row) at (3,\y) {$b_2$};
  \node[girl] (g2\row) at (4,\y) {$g_2$};
  \node[boy,free] (b3\row) at (5,\y) {$b_3$};
}
\draw[edge,line width=.8pt,draw=black!60] (g00)--(b10) (g10)--(b20) (g20)--(b30);
\draw[couple] (b10)--(g10) (b20)--(g20);
\draw[couple] (g01)--(b11) (g11)--(b21) (g21)--(b31);
\draw[edge,line width=.8pt,draw=black!60] (b11)--(g11) (b21)--(g21);
\node[girl] at (g01) {$g_0$};\node[boy] at (b31) {$b_3$};
\node[anchor=east,font=\small] at (-.5,0) {augmenting path};
\node[anchor=east,font=\small] at (-.5,-1.25) {after augmenting};
\draw[-{Stealth},black!50] (2.5,-.38)--(2.5,-.87);
\end{tikzpicture}

\caption{An augmenting path of level~$L=3$, and the matching after
augmenting along it. Double lines are couples; heavy outlines mark free
vertices. In all figures girls are drawn as circles and boys as squares.}
\label{fig:augpath}
\end{figure}

Berge's observation~\cite{Berge} (pre-fascicle, page~1) is that if $M'$ has
$s'$ couples, then the union of $M'$ with a maximum matching contains at
least $s-s'$ \emph{vertex-disjoint} augmenting paths for $M'$. We shall
use this mostly when $s=M=N$, so that there is a \emph{perfect}
matching~$P$. Then the symmetric difference $M'\symdiff P$ consists of
exactly $f$ vertex-disjoint augmenting paths, where $f=M-s'$ is the number
of free girls, together with some alternating cycles. Following answer
15(a), we call these $f$ paths the \emph{basic paths} of $M'$ with
respect to~$P$. A basic path of level~$\ell$ contains exactly $\ell$ edges
of~$P$, and like Knuth in answer~15(d) we often draw it as a \emph{row} of
$\ell$ cells, one for each girl--boy pair of~$P$ that it visits. (Later
sections call these rows the rows of a matrix, or hooks; they are always
basic paths.)

\subsection{Shortest augmenting paths and rounds}
\label{sec:dag}

A \emph{shortest augmenting path}, or SAP, is an augmenting path of
smallest level. To find them, orient every edge relative to the current
matching: an edge outside the matching points from its girl to its boy,
and a couple points from its boy to its girl. Alternating paths that start
at a free girl are then exactly the directed paths that start there. The
\emph{distance} of a vertex is the number of edges on a shortest such
directed path from some free girl to it.

Hopcroft and Karp~\cite{HK} organize all SAPs into a directed acyclic
graph, built breadth-first from the free girls (page~2 of the
pre-fascicle). The dag has \emph{levels} $G_0,B_1,G_1,\ldots,
G_{L-1},B_L$: $G_0$ is the set of free girls; $B_k$ consists of the
boys first reached from girls in $G_{k-1}$; and $G_k$ consists of the
mates of the boys in $B_k$. So $G_k$ and $B_k$ are the girls at distance
$2k$ and the boys at distance $2k-1$. There is an arc $b\to g$ for each edge
$g\sg b$ with $g\in G_{k-1}$, $b\in B_k$, and an arc from each girl to
her mate. The construction stops at the first level~$B_L$ that contains
a free boy. Knuth adds two special nodes: arcs $g\to\bot$ from the free
girls and $\top\to b$ to the free boys of~$B_L$. The SAPs are then
exactly the paths from $\top$ to $\bot$, read backwards; this is why the
arcs in the pre-fascicle's figure~(7) point to the left.

Each \emph{round} of Algorithm~H builds this dag (steps~\Hs2--\Hs{10})
and then dismantles it (steps~\Hs{11}--\Hs{17}): it repeatedly finds a
path from $\top$ to $\bot$ by depth-first search, augments the matching
along it, and deletes its vertices from the dag. When no path remains,
the paths removed form a \emph{maximal} set of vertex-disjoint SAPs:
no further SAP is disjoint from all of them. This set need not be as
large as possible, and the gap between ``maximal'' and ``maximum'' is
what a worst-case input exploits.

After a round at level~$L$, every augmenting path has level greater
than~$L$; this is exercise~5, and goes back to Hopcroft and Karp
\cite[Theorem~2 and Corollaries~3--4]{HK}. So the levels of successive rounds strictly
increase. Theorem~H turns this into a
counting bound: when a round at level~$L$ is over, every one of the
$s-s'$ disjoint augmenting paths has at least $L+1$ edges of a maximum
matching, so $s\ge(L+1)(s-s')$. Corollary~K deduces that there are at
most about $2\sqrt s$ rounds, each costing $O(M+N+E)$.

\subsection{What an input is}
\label{sec:input}

Algorithm~H is deterministic. To speak of a graph that ``might cause''
it to behave badly, as exercise~15(c) does, we must say exactly which
features of the input it reads. There are two:
\begin{enumerate}[label=(\roman*)]
\item the numbering of the girls~$1,\ldots,M$, which fixes the order in
  which step~\Hs1 scans them;
\item for each girl~$g$, the order of her list of boys, which starts at
  $\A{GLINK}[g]$ and continues through $\A{GNEXT}$.
\end{enumerate}
We call a bipartite graph together with these data an \emph{ordered
graph}. By exercise~6, any desired list order arises from a suitable
order of the input pairs, since each new edge is inserted at the front
of its girl's list. Nothing else matters. The numbering of the boys is
irrelevant, since no step compares boy indices. There is also \emph{no}
boy-side order to choose: the reverse lists $\A{BLINK}[b]$ are built by
the algorithm itself.

For orientation, here is what the steps do, in our words (the precise
code is on pages~4--5 of the pre-fascicle).
\begin{center}\small
\begin{tabular}{@{}>{\raggedright\arraybackslash}p{.17\textwidth}>{\raggedright\arraybackslash}p{.76\textwidth}@{}}
\toprule
\Hs1&Greedy start: a maximal matching; the free girls go into \A{QUEUE}.\\
\Hs2--\Hs3&Start a dag; start the next level.\\
\Hs4--\Hs7&Take the next girl of the current level and scan her list of boys.\\
\Hs8--\Hs9&A boy seen for the first time joins the next level; his mate is
  queued for the level after, or, if he is free, the final level~$L$ is
  fixed and he is recorded as a root.\\
\Hs{10}&Go on to the next level, or finish the dag (or stop if no free boy was reached).\\
\Hs{11}&Take the next root; if none is left, reset all marks and start a new round.\\
\Hs{12}--\Hs{15}&Depth-first search from the root down the reverse lists
  \A{BLINK}, consuming each arc it follows and skipping used girls;
  back up when a list runs out.\\
\Hs{16}--\Hs{17}&A free girl was reached: remove her from the free list and
  augment along the path.\\
\bottomrule
\end{tabular}
\end{center}

The following consequences of the code will be used repeatedly.

\begin{itemize}
\item \textbf{\Hs1 scans the girls in decreasing order.} Answer~7 matches
  girl~$g$ to the first free boy on her list, for $g=M,M-1,\ldots,1$,
  and appends each unsuccessful girl to \A{QUEUE}. So the free girls
  start out in \A{QUEUE} in decreasing numerical order. The resulting
  matching is maximal, so the first dag has $L\ge2$.
\item \textbf{The dag levels are scanned in queue order.} Girls of $G_l$
  are processed in their order in \A{QUEUE}; each scans her own list in
  order. A boy enters $B_{l+1}$ when he is first seen (step~\Hs8), and
  his mate is appended to \A{QUEUE} as a girl of $G_{l+1}$ (step
  \Hs9). So the order of $G_{l+1}$ is the order in which the boys of
  $B_{l+1}$ were discovered.
\item \textbf{Reverse lists are built by insertion at the head.} The
  first girl who sees $b$ becomes the only entry of $\A{BLINK}[b]$
  (step~\Hs8); every later girl of the same level who is adjacent to $b$
  is inserted at the \emph{front} (step~\Hs6). Hence $\A{BLINK}[b]$ lists
  the girls of $G_l$ adjacent to $b$ in the \emph{reverse} of their order
  in \A{QUEUE}.
\item \textbf{The last level is cut short.} When the first free boy is
  found, $L$ is set to the current level and the queue of the next
  level is discarded. After that, matched boys are no longer entered
  (step~\Hs8), but free boys still are. The free boys of $B_L$ are the
  last $r$ entries of \A{MARKED}, in order of discovery, and step~\Hs{11}
  uses them as roots of the depth-first search in that order.
\item \textbf{The depth-first search destroys what it reads.} Step~\Hs{13}
  removes each arc from $\A{BLINK}[b]$ as it follows it, and skips a
  girl if she has already been used by an earlier SAP of this round
  (she is ``deleted''). If a list becomes empty, step~\Hs{15} backs up
  one level; nothing is ever re-examined. Reaching a free girl completes
  a SAP.
\item \textbf{A used free girl is replaced by the last one.} Step~\Hs{16}
  removes the girl~$g_0$ of a new SAP from the free list by moving the
  last free girl, $\A{QUEUE}[f-1]$, into her slot. So the order of the
  free girls is preserved only if each SAP ends at the last of them.
\end{itemize}

\subsection{A worked round}
\label{sec:worked}

The ordered graph of Fig.~\ref{fig:gamma5} has nine girls and nine
boys. Its girl lists, in increasing order of girls, are
\begin{equation}
  [6,8];\ [4,6];\ [1,4];\ [1];\ [1,2];\ [2,3];\ [4,2,5];\ [6,5,7];\ [8,9].
  \label{eq:gamma5lists}
\end{equation}
(Girl~1 is adjacent to boys~6 and~8, in that order, and so on.)
Its vertices are arranged in nine \emph{cells}: cell~$(a,c)$, in row~$a$
and column~$c$, holds one girl and one boy, joined by an edge (the short
dark edge inside the cell). These nine edges form a perfect matching, so
$s=9$. The other edges join the boy of a cell either to the girl of the
next cell in the same row or to the girl of the cell below; in the
drawing they run diagonally, up-right and down-left respectively. We
shall meet this graph again in Section~\ref{sec:staircase} as the
staircase graph~$\Gam_5$.

\begin{figure}[ht]
\centering
\begin{tikzpicture}[]
\begin{scope}[on background layer]
\end{scope}
\draw[edge] (-0.468,-5.967) -- (0.468,-4.563);
\draw[pedge] (-0.468,-5.967) -- (0.468,-6.669);
\draw[edge] (-0.468,-3.861) -- (0.468,-2.457);
\draw[pedge] (-0.468,-3.861) -- (0.468,-4.563);
\draw[edge] (-0.468,-1.755) -- (0.468,-0.351);
\draw[pedge] (-0.468,-1.755) -- (0.468,-2.457);
\draw[pedge] (-0.468,0.351) -- (0.468,-0.351);
\draw[couple] (2.002,0.351) -- (0.468,-0.351);
\draw[pedge] (2.002,0.351) -- (2.938,-0.351);
\draw[couple] (4.472,0.351) -- (2.938,-0.351);
\draw[pedge] (4.472,0.351) -- (5.408,-0.351);
\draw[couple] (2.002,-1.755) -- (0.468,-2.457);
\draw[edge] (2.002,-1.755) -- (2.938,-0.351);
\draw[pedge] (2.002,-1.755) -- (2.938,-2.457);
\draw[couple] (2.002,-3.861) -- (0.468,-4.563);
\draw[edge] (2.002,-3.861) -- (2.938,-2.457);
\draw[pedge] (2.002,-3.861) -- (2.938,-4.563);
\draw[couple] (2.002,-5.967) -- (0.468,-6.669);
\draw[pedge] (2.002,-5.967) -- (2.938,-6.669);
\node[girl,free] at (-0.468,-5.967) {1};
\node[girl,free] at (-0.468,-3.861) {2};
\node[girl,free] at (-0.468,-1.755) {3};
\node[girl,free] at (-0.468,0.351) {4};
\node[girl] at (2.002,0.351) {5};
\node[girl] at (4.472,0.351) {6};
\node[girl] at (2.002,-1.755) {7};
\node[girl] at (2.002,-3.861) {8};
\node[girl] at (2.002,-5.967) {9};
\node[boy] at (0.468,-0.351) {1};
\node[boy] at (2.938,-0.351) {2};
\node[boy,free] at (5.408,-0.351) {3};
\node[boy] at (0.468,-2.457) {4};
\node[boy,free] at (2.938,-2.457) {5};
\node[boy] at (0.468,-4.563) {6};
\node[boy,free] at (2.938,-4.563) {7};
\node[boy] at (0.468,-6.669) {8};
\node[boy,free] at (2.938,-6.669) {9};
\node[lab,font=\scriptsize,text=black!50,anchor=west] at (-0.078,0.468) {$(1,1)$};
\node[lab,font=\scriptsize,text=black!50,anchor=west] at (2.392,0.468) {$(1,2)$};
\node[lab,font=\scriptsize,text=black!50,anchor=west] at (4.862,0.468) {$(1,3)$};
\node[lab,font=\scriptsize,text=black!50,anchor=west] at (-0.078,-1.638) {$(2,1)$};
\node[lab,font=\scriptsize,text=black!50,anchor=west] at (2.392,-1.638) {$(2,2)$};
\node[lab,font=\scriptsize,text=black!50,anchor=west] at (-0.078,-3.744) {$(3,1)$};
\node[lab,font=\scriptsize,text=black!50,anchor=west] at (2.392,-3.744) {$(3,2)$};
\node[lab,font=\scriptsize,text=black!50,anchor=west] at (-0.078,-5.850) {$(4,1)$};
\node[lab,font=\scriptsize,text=black!50,anchor=west] at (2.392,-5.850) {$(4,2)$};
\draw[edgecol,line width=.5pt,dash pattern=on 1.2pt off 1.6pt] (2.938,-4.563) -- (2.002,-5.967);
\node[lab,anchor=west,font=\scriptsize] at (2.610,-5.315) {omitted};
\end{tikzpicture}

\caption{The 9-pair example, with the matching produced by~\Hs1. The
grey label $(a,c)$ names the cell of the girl beside it and of the boy
below her; the dark edge joining them belongs to a perfect matching.
Light grey lines are the other edges of the graph, double lines are
couples, and free vertices have heavy outlines. The dotted segment
marks an edge deliberately left out of the graph
(Section~\ref{sec:gamma}).}
\label{fig:gamma5}
\end{figure}

\paragraph{Step H1.}
Girls~$9,8,7,6,5$ come first. Each takes the first boy on her list,
namely $8,6,4,2,1$, and these are all different. Girls~$4,3,2,1$ then
find every boy on their lists taken. So~\Hs1 produces the five couples
drawn as double lines in Fig.~\ref{fig:gamma5}, and
\[
   \A{QUEUE}=(4,3,2,1),\qquad f=4 .
\]
In the picture, the free girls are the first girls of the four rows,
and the free boys are the last boys of the rows.

\paragraph{Building the dag.}
Fig.~\ref{fig:gamma5dag} shows the dag of the first round. Girl~4 scans
her list $[1]$ and discovers boy~1. Girl~3 scans $[1,4]$: boy~1 has
already been seen at this level, so girl~3 is inserted at the head of
$\A{BLINK}[1]$, giving $\A{BLINK}[1]=(3,4)$; then she discovers boy~4.
Girls~2 and~1 behave the same way, so $B_1=\{1,4,6,8\}$ in that order,
and their mates $5,7,8,9$ form $G_1$. At the next level boy~2 is
discovered first; he is matched (to girl~6), so the construction goes
on. Then girl~7 discovers boy~5, who is free: now $L=2$. Boys~7 and~9
are free too, so the roots of the search will be $5,7,9$, in that
order.

\begin{figure}[ht]
\centering
\begin{tikzpicture}[]
\begin{scope}[on background layer]
\draw[sap] (0.000,-1.230) -- (1.550,-2.460) -- (3.100,-1.640) -- (4.650,-1.640) -- (6.200,-0.820) -- (7.750,-1.230);
\end{scope}
\node[girl] (g4) at (1.550,0.000) {$g_{4}$};
\node[girl] (g3) at (1.550,-0.820) {$g_{3}$};
\node[girl] (g2) at (1.550,-1.640) {$g_{2}$};
\node[girl] (g1) at (1.550,-2.460) {$g_{1}$};
\node[boy] (b1) at (3.100,0.000) {$b_{1}$};
\node[boy] (b4) at (3.100,-0.820) {$b_{4}$};
\node[boy] (b6) at (3.100,-1.640) {$b_{6}$};
\node[boy] (b8) at (3.100,-2.460) {$b_{8}$};
\node[girl] (g5) at (4.650,0.000) {$g_{5}$};
\node[girl] (g7) at (4.650,-0.820) {$g_{7}$};
\node[girl] (g8) at (4.650,-1.640) {$g_{8}$};
\node[girl] (g9) at (4.650,-2.460) {$g_{9}$};
\node[boy] (b2) at (6.200,0.000) {$b_{2}$};
\node[boy] (b5) at (6.200,-0.820) {$b_{5}$};
\node[boy] (b7) at (6.200,-1.640) {$b_{7}$};
\node[boy] (b9) at (6.200,-2.460) {$b_{9}$};
\node[circle,draw,inner sep=1pt,minimum size=5mm] (bot) at (0.000,-1.230) {$\bot$};
\node[circle,draw,inner sep=1pt,minimum size=5mm] (top) at (7.750,-1.230) {$\top$};
\node[girl,dashed,opacity=.45] (gx6) at (7.200,0.550) {$g_{6}$};
\draw[arc,dashed,opacity=.45] (gx6) -- (b2);
\draw[arc] (g4) -- (bot);
\draw[arc] (g3) -- (bot);
\draw[arc] (g2) -- (bot);
\draw[arc] (g1) -- (bot);
\draw[arc] (g5) -- (b1);
\draw[arc] (g7) -- (b4);
\draw[arc] (g8) -- (b6);
\draw[arc] (g9) -- (b8);
\draw[arc] (b1) -- node[pos=.24,fill=white,inner sep=.5pt,font=\scriptsize] {1} (g3);
\draw[arc] (b1) -- node[pos=.24,fill=white,inner sep=.5pt,font=\scriptsize] {2} (g4);
\draw[arc] (b4) -- node[pos=.24,fill=white,inner sep=.5pt,font=\scriptsize] {1} (g2);
\draw[arc] (b4) -- node[pos=.24,fill=white,inner sep=.5pt,font=\scriptsize] {2} (g3);
\draw[arc] (b6) -- node[pos=.24,fill=white,inner sep=.5pt,font=\scriptsize] {1} (g1);
\draw[arc] (b6) -- node[pos=.24,fill=white,inner sep=.5pt,font=\scriptsize] {2} (g2);
\draw[arc] (b8) --  (g1);
\draw[arc] (b2) -- node[pos=.24,fill=white,inner sep=.5pt,font=\scriptsize] {1} (g7);
\draw[arc] (b2) -- node[pos=.24,fill=white,inner sep=.5pt,font=\scriptsize] {2} (g5);
\draw[arc] (b5) -- node[pos=.24,fill=white,inner sep=.5pt,font=\scriptsize] {1} (g8);
\draw[arc] (b5) -- node[pos=.24,fill=white,inner sep=.5pt,font=\scriptsize] {2} (g7);
\draw[arc] (b7) --  (g8);
\draw[arc] (b9) --  (g9);
\draw[arc] (top) -- (b5);
\draw[arc] (top) -- (b7);
\draw[arc] (top) -- (b9);
\node at (1.550,-3.210) {$G_0$};
\node at (3.100,-3.210) {$B_1$};
\node at (4.650,-3.210) {$G_1$};
\node at (6.200,-3.210) {$B_2$};
\end{tikzpicture}

\caption{The first dag for the graph of Fig.~\ref{fig:gamma5}, drawn as
in the pre-fascicle's figure~(7). When a boy has two arcs, the small
numbers give their order in $\A{BLINK}$; the girl who discovered him
comes \emph{second}. The shaded path is the only SAP that Algorithm~H
finds. Boy~2 belongs to $B_2$ but is matched, so his mate (faded) is
not part of the dag.}
\label{fig:gamma5dag}
\end{figure}

\paragraph{Dismantling the dag.}
The search starts at boy~5. The head of his list is girl~8, whose mate
is boy~6; the head of $\A{BLINK}[6]$ is girl~1, who is free. That
completes the SAP
\[
    g_1\sg b_6\db g_8\sg b_5 ,
\]
and Algorithm~H augments along it. The next root, boy~7, has only
girl~8 on his list, and she has just been used. The last root, boy~9,
leads to girl~9 and then to boy~8, whose only remaining entry is
girl~1, also used. So the round ends with a \emph{single} SAP, even
though the four rows of Fig.~\ref{fig:gamma5} contain three disjoint
SAPs (rows~2, 3, and~4). The one path that was found meets all three
of them. Step~\Hs{16} removed girl~1, the last free girl, so the queue
for the next round is $(4,3,2)$.

The remaining rounds behave similarly: each finds one SAP, at levels
$3$, $4$, and $5$, so the whole run has five rounds counting~\Hs1.
Fig.~\ref{fig:gamma5rounds} on page~\pageref{fig:gamma5rounds} shows
all of them. A Tchoukaillon graph with $s=9$ would need only four.

\subsection{Counting rounds and operations}
\label{sec:counting}

Answer~15(c) gives the number of rounds of a Tchoukaillon graph as
``the final value of~$L$.'' In such a graph the dags have levels
$2,3,\ldots,L$, so this count treats step~\Hs1 as a round at level~1;
answer~7 makes the same identification when it notes that, if~\Hs1 did
not produce a maximal matching, the first round would do so with $L=1$.
We adopt this convention.

\begin{definition}
A round is \emph{successful} if it enlarges the matching. The
\emph{number of rounds} of a run of Algorithm~H is the number of
successful rounds, where step~\Hs1 counts as one round (successful
whenever the graph has an edge). The final dag construction, which
finds no augmenting path and stops the algorithm, is not counted.
\end{definition}

Thus a run on a graph with $s\ge1$ has $1+k$ rounds if $k$ dags contain
augmenting paths. If the levels of those dags are $2,3,\ldots,n$ then
there are exactly $n$ rounds. The run of Section~\ref{sec:worked} has
five.

For running times we count elementary operations of Algorithm~H rather
than seconds. In each successful round we record three counts, named
after the steps that perform them:
\begin{align*}
  W_5&=\text{number of adjacency entries examined in step~\Hs5 (the breadth-first search)},\\
  W_{13}&=\text{number of arcs consumed in step~\Hs{13} (the depth-first search)},\\
  W_{11}&=\text{number of marks reset in step~\Hs{11}, which equals the number of boys in the dag.}
\end{align*}
$W_{13}$ includes arcs that lead to deleted girls. Every other step of a
round costs $O(1)$ per unit of one of these counts, plus $O(f)$, so the
total work in rounds after~\Hs1 is $\Theta(W_5+W_{13}+W_{11})$ up to that
small additive term. For example, the round of Section~\ref{sec:worked}
has $W_5=17$ (girls~$4,3,2,1,5,7,8,9$ have lists of total length~17),
$W_{11}=8$ (the boys of $B_1$ and $B_2$), and $W_{13}=5$ (two arcs for the
successful root, one for boy~7, and two for boy~9).
